\section{Models}
\subsection{Quantum neural network}
The QNN is a parameterized variational circuit with data re-uploading. For two-dimensional datasets, it operates on a two-qubit register ($n=2$); for the four-dimensional QDataSet dynamics benchmark, it scales to an $n=4$ qubit register. Within each layer $\ell$, single-qubit rotations $R_y(x_i)$ embed each normalized input feature into qubit $i$, followed by trainable rotations $R_y(\theta_{\ell,i})$. In entangled configurations, CNOT gates couple adjacent qubits (with qubit $0\to1$ for 2 qubits, and a circular ladder for 4 qubits). This layer repeats $L$ times (primary $L=2$; depths 1 and 3 are also evaluated).

For the 2-qubit $\ell$th layer, the unitary operation is
\begin{equation}
U_\ell(x,\theta_\ell)=\operatorname{CNOT}_{0\rightarrow1}
\left[R_y(\theta_{\ell,0})\otimes R_y(\theta_{\ell,1})\right]
\left[R_y(x_0)\otimes R_y(x_1)\right],
\end{equation}
where the CNOT is omitted for the unentangled ablation. The final state is $|\psi(x,\theta)\rangle=U_L\cdots U_1|0\rangle^{\otimes n}$. The quantum circuit returns all single-qubit expectation values $h(x,\theta)=[\langle Z_0\rangle,\dots,\langle Z_{n-1}\rangle]^\top$. A trainable classical linear readout combines them into a logit $s=w^\top h+b$, which is converted to a positive-class probability $p(y=1\mid x)=\sigma(s)$. For a labelled training set $\{(x_i,y_i)\}_{i=1}^N$, the objective is binary cross-entropy:
\begin{equation}
  \mathcal{L}(w,b,\theta)=-\frac{1}{N}\sum_{i=1}^{N}\left[y_i\log p_i+(1-y_i)\log(1-p_i)\right].
\end{equation}
This establishes a hybrid quantum-classical pipeline: the PennyLane QNode computes quantum expectation values, PyTorch computes the linear readout and loss, and automatic differentiation propagates gradients through both components. The primary implementation executes on PennyLane's \texttt{default.qubit} device using analytic simulator backpropagation, verified against an independent PyTorch statevector reference.

\begin{figure}[!htbp]
\centering
\begin{tikzpicture}[x=1.35cm,y=1.28cm,font=\normalsize]
  % Qubit wires
  \draw[wire,line width=1.1pt] (0,1) -- (9.45,1);
  \draw[wire,line width=1.1pt] (0,0) -- (9.45,0);
  \node[anchor=east,font=\large\bfseries] at (-0.12,1) {$q_0:|0\rangle$};
  \node[anchor=east,font=\large\bfseries] at (-0.12,0) {$q_1:|0\rangle$};
  % Layer 1 encoding
  \node[gate,fill=qblue,minimum width=1.35cm,minimum height=0.68cm,font=\normalsize] at (0.95,1) {$R_y(x_0)$};
  \node[gate,fill=qblue,minimum width=1.35cm,minimum height=0.68cm,font=\normalsize] at (0.95,0) {$R_y(x_1)$};
  \node[gate,fill=qgold,minimum width=1.35cm,minimum height=0.68cm,font=\normalsize] at (2.42,1) {$R_y(\theta_{10})$};
  \node[gate,fill=qgold,minimum width=1.35cm,minimum height=0.68cm,font=\normalsize] at (2.42,0) {$R_y(\theta_{11})$};
  \fill[qink] (3.52,1) circle (2.6pt);
  \draw[qink,line width=1.2pt] (3.52,1) -- (3.52,0);
  \draw[qink,line width=1.2pt,fill=white] (3.52,0) circle (5.5pt);
  \draw[qink,line width=1.2pt] (3.52,-0.10) -- (3.52,0.10);
  \draw[qink,line width=1.2pt] (3.42,0) -- (3.62,0);
  % Layer 2 encoding
  \node[gate,fill=qblue,minimum width=1.35cm,minimum height=0.68cm,font=\normalsize] at (4.55,1) {$R_y(x_0)$};
  \node[gate,fill=qblue,minimum width=1.35cm,minimum height=0.68cm,font=\normalsize] at (4.55,0) {$R_y(x_1)$};
  \node[gate,fill=qgold,minimum width=1.35cm,minimum height=0.68cm,font=\normalsize] at (6.02,1) {$R_y(\theta_{20})$};
  \node[gate,fill=qgold,minimum width=1.35cm,minimum height=0.68cm,font=\normalsize] at (6.02,0) {$R_y(\theta_{21})$};
  \fill[qink] (7.12,1) circle (2.6pt);
  \draw[qink,line width=1.2pt] (7.12,1) -- (7.12,0);
  \draw[qink,line width=1.2pt,fill=white] (7.12,0) circle (5.5pt);
  \draw[qink,line width=1.2pt] (7.12,-0.10) -- (7.12,0.10);
  \draw[qink,line width=1.2pt] (7.02,0) -- (7.22,0);
  % Measurements
  \node[draw=qline,rounded corners=2pt,fill=qteal,line width=1.0pt,minimum width=1.12cm,minimum height=0.66cm,font=\normalsize\bfseries] at (8.25,1) {$\langle Z_0\rangle$};
  \node[draw=qline,rounded corners=2pt,fill=qteal,line width=1.0pt,minimum width=1.12cm,minimum height=0.66cm,font=\normalsize\bfseries] at (8.25,0) {$\langle Z_1\rangle$};
  \node[draw=qline,rounded corners=3pt,fill=qcoral,line width=1.0pt,minimum width=1.65cm,minimum height=1.15cm,align=center,font=\small\bfseries] (readout) at (9.92,0.5) {Linear\\readout\\+ sigmoid};
  \draw[arrow,line width=1.2pt] (8.68,1) -- (readout.west |- 0,0.74);
  \draw[arrow,line width=1.2pt] (8.68,0) -- (readout.west |- 0,0.26);
  \node[font=\small\bfseries,align=center] at (2.4,1.60) {Layer 1};
  \node[font=\small\bfseries,align=center] at (6.0,1.60) {Layer 2};
\end{tikzpicture}
\caption{Two-qubit circuit used by the primary QNN. Each layer re-encodes both features. The CNOT control is on $q_0$ and the target is on $q_1$. The depth-1 and depth-3 variants add or remove a full layer.}
\label{fig:circuit}
\end{figure}

\subsection{Classical baseline}
The wide classical model is a PyTorch MLP with architecture $d\rightarrow16\rightarrow8\rightarrow1$, ReLU activations between hidden layers, and a single binary logit at the output. It contains 193 trainable parameters on 2D datasets and 225 parameters on the 4D QDataSet dynamics benchmark. To establish rigorous capacity-controlled baselines, we also evaluate compact $d\rightarrow2\rightarrow1$ MLPs with 9 parameters ($d=2$) and 13 parameters ($d=4$). For depth-2 QNNs, the 2-qubit circuit contains 4 gate angles and 3 readout parameters (7 parameters total), while the 4-qubit circuit contains 8 gate angles and 5 readout parameters (13 parameters total), establishing an exact 1:1 parameter match (13 vs. 13) on the QDataSet benchmark.

\subsection{Optimization}
Both models use full-batch Adam optimization and binary cross-entropy with logits for 100 epochs. The MLP learning rate is 0.01; the QNN learning rate is 0.025. Weight decay is $10^{-4}$ and gradients are clipped to a maximum norm of 5. The different fixed learning rates reflect the different parameterizations; no exhaustive hyperparameter search was conducted. The primary experiment uses exact analytic expectation values during training.
